\scanpage{159}
% 159.01 Top centre
\mspage{487}

% 159.02 Upper paragraph, continuation
rendre compte en passant, quand je me voyais confronté à nouveau à ce sempiternel mécanisme en moi de r\,e\,j\,e\,t de tout ce qui se présente sous une apparence déplaisante, chez autrui ou chez moi-même. (Mais quand il d'agit de soi-même, ce mécanisme le plus souvent a comme effet de ne pas même prendre connaissance de la chose déplaisante en question...)

% 159.03 Main paragraph
L'acceptation dont je parle prend racine dans un i\,n\,t\,é\,r\,ê\,t pour cette chose qu'on ``accepte'', chez soi-même ou chez autrui. Alors que l'acceptation est par elle-même une disposition intérieure de caractère typiquement ``yin''\unsure{:|,} cette connot\handwritten{a}tion d'``intérêt'' qu'elle prend chez moi est de nature ``yang'' - \unsure{c'es|c'est} le ``yang dans le yin'', dans la délicate dialectique chinoise de l'entrelacement à l'infini du yin et du yang... J'allais me hasarder à dire, un peu dans la foulée\handwritten{,} qu'il y avait une identité pure et simple entre l'acceptation (la vraie!) et cet intérêt, cette curiosité. Pourtant, en posant un peu sur la chose, je me rends compte qu'il y a bien aussi une autre façon d'accepter, de nature, elle, plus totalement yin que celle qui m'est surtout coutumière. C'est comme un a\,c\,c\,u\,e\,i\,l de la chose acceptée, et non un élan vers elle pour la sonder. (Cette nuance d'accueil m'apparaît du coup comme le ``y\,i\,n d\,ans le yin'', nous y voilà!) L'élan de l'intérêt, et l'attitude d'accueil, peuvent l'une comme l'autre former la note de fond de l'acceptation d'autrui ou de soi. La chose commune aux deux est la s\,y\,m\,p\,a\,t\,h\,i\,e . C'est là aussi une des formes de l' a\,m\,o\,u\,r . S'il y a quelque identité profonde à cerner ici, ce serait donc par la constatation que l\,'\,a\,c\,c\,e\,p\,t\,a\,t\,i\,o\,n e\,s\,t u\,n\,e f\,o\,r\,m\,e d\,e l\,'\,a\,m\,o\,u\,r . Amour de soi, amour de l'autre, l'un et l'autre indissolublement liés...

% 159.04 Lower paragraph, ending mid-sentence
Sauf en de rares instants, mon intérêt est plus intensément impliqué quand il s'agit de ma propre personne, que de celle d'autrui. C'est cet intérêt passioné pour ma personne qui a animé les longues périodes de méditation, au cours de ces huit dernières années. Il est vrai que c'est bien la connaissance de soi qui est au coeur de la connaissance d'autrui et du monde, et non l'inverse - et je sens que c'est bien vers le coeur des choses, vers le plus essentiel, que m'a porté et me porte encore ma nouvelle passion, la méditation. L'intérêt pour autrui est apparu de façon plus parcellaire et plus réticente au cours de ces années, tout comme l'acceptation qui en découle. Une des façons dont elle s'est manifestée concrètement, c'est par une propension moins grande à parler quand je suis en compagnie, et par une attitude d'écoute. La plus grande partie de ma vie, cette capacité d'écoute m'avait presque entièrement fait défaut. Encore après le grand tournant des retrouvailles, il m'a fallu constater bien souvent encore
% archive sheet number: 158
